\section{Peripheral spectral refinements}

Every result in this section concerning peripheral eigenvalues, cyclic
projections, or peripheral eigenspaces is stated for a finite Kraus map and
therefore assumes complete positivity. These results do not supply the abstract
Schwarz-map case of Wolf Theorem~6.6; the restriction and its elimination plan
are documented in the scope-restriction note
\cite{gap:wolf_thm6_6_kraus_scope}.


\subsection{Cyclic decomposition of irreducible finite Kraus maps}

The construction --- normalizing a peripheral eigenvalue to a unitary
eigenvector, forming its discrete Fourier (cyclic) spectral projections,
and the supporting corner-algebra apparatus (restriction to an invariant
corner, corner rank, and the compression isometry
$P\MN{D}P\cong\MN{n}$) --- is given in
Section~\ref{sec:app_spectral_periodicity}.


\begin{lemma}[Cyclic shift of the Kraus operators]
    \label{lem:kraus_cyclic_shift}
    \lean{Kraus.one_sub_mul_kraus_mul_eq_zero_of_mapLM_proj}
    \lean{Kraus.kraus_mul_cyclicProj}
    \lean{Kraus.cyclicProj_ne_zero}
    \leanok
    Let $P_0,\ldots,P_{m-1}$ be orthogonal projections summing to the
    identity and cyclically permuted by a Kraus map,
    $\sum_vK_vP_{k+1}K_v^\dagger=P_k$. Then the projections are mutually
    orthogonal, none of them vanishes, and each Kraus operator intertwines
    consecutive projections:
    \begin{align}
        K_vP_{k+1}
         & =P_kK_v.
        \label{eq:spectral_kraus_cyclic_shift}
    \end{align}
\end{lemma}

\begin{proof}
    \leanok

    Compressing $\sum_vK_vP_{k+1}K_v^\dagger=P_k$ by $\Id-P_k$ exhibits a
    vanishing sum of positive semidefinite matrices, so each summand
    vanishes: $(\Id-P_k)K_vP_{k+1}=0$. Mutual orthogonality follows by
    compressing the resolution of the identity by a fixed $P_k$, another
    vanishing sum of positive semidefinite matrices. Expanding
    $K_vP_{k+1}$ and $P_kK_v$ through the resolution of the identity, the
    off-diagonal compressions vanish by orthogonality, leaving
    (\ref{eq:spectral_kraus_cyclic_shift}). If some projection vanished,
    the cyclic action would propagate the vanishing around the cycle,
    contradicting the resolution of the identity.
\end{proof}

\begin{theorem}[Cyclic shift along a Kraus word]
    \label{thm:eval_word_cyclic_shift}
    \lean{MPSTensor.evalWord_mul_cyclicProj}
    \leanok
    \uses{lem:kraus_cyclic_shift}
    Under the hypotheses of Lemma~\ref{lem:kraus_cyclic_shift}, a product of
    $\ell$ Kraus operators moves each cyclic projection back by $\ell$ steps.
    In particular, every product of $m$ Kraus operators commutes with every
    cyclic projection.
\end{theorem}

\begin{proof}\leanok
    \uses{lem:kraus_cyclic_shift}
    Induct on the length of the word using
    (\ref{eq:spectral_kraus_cyclic_shift}).
\end{proof}

The primitivity and irreducibility of the restricted sector dynamics, the
invariance of each cyclic corner under the period, and the closure of the
peripheral eigenvalues under products and inverses are proved in
Section~\ref{sec:app_spectral_periodicity}.


\section{Self-overlap and the period-one criterion}

\begin{theorem}[Self-overlap along peripheral-period multiples]
    \label{thm:bnt_periodic_self_overlap_tendsto}
    \label{thm:periodic_self_overlap_tendsto}
    \lean{MPSTensor.periodicSelfOverlap_tendsto}
    \leanok
    \uses{def:mpv_overlap, def:irreducible_tensor, def:left_canonical_tensor, def:blocked_tensor}
    Let $D>0$ and $m>0$. Let $A$ be a left-canonical irreducible MPS tensor
    with bond dimension $D$, and suppose that the peripheral eigenvalues of
    its transfer map are exactly $\{z\in\C:z^m=1\}$. Then
    $O_{AA}(mn)\longrightarrow m$.
\end{theorem}

\begin{proof}
    \leanok
    \uses{thm:cyclic_sector_decomp_after_blocking_isometry, thm:blocked_sector_unconditional, thm:primitive_overlap, thm:overlap_decay_irr_tp, thm:overlap_decay_rect_irr_tp}
    Apply
    Theorem~\ref{thm:cyclic_sector_decomp_after_blocking_isometry} to the
    period-$m$ blocking $A^{[m]}$. It gives $m$ nonzero cyclic-sector
    tensors $C_0,\ldots,C_{m-1}$, each left-canonical, whose unit-weight direct
    sum has the same MPV family as $A^{[m]}$. The cyclic-sector relations make
    every $C_u$ primitive and irreducible. Distinct sectors are not
    gauge-phase equivalent. Hence the primitive self-overlap limit and the
    equal- and unequal-dimension overlap-decay theorems give
    \begin{align}
        O_{C_uC_v}(n)
         & \longrightarrow
        \begin{cases}
            1, & u=v,     \\
            0, & u\neq v.
        \end{cases}
        \notag
    \end{align}
    Expanding the overlap of the unit-weight direct sum and taking the finite
    sum of these limits yields
    \begin{align}
        O_{A^{[m]}A^{[m]}}(n)
         & =\sum_{u,v\in\Z_m}O_{C_uC_v}(n)
        \longrightarrow m.
        \notag
    \end{align}
    Finally, blocking identifies
    $O_{A^{[m]}A^{[m]}}(n)=O_{AA}(mn)$, which proves the claim.
\end{proof}

\begin{theorem}[Primitivity and normality from normalized self-overlap]
    \label{thm:bnt_primitivity_normality_from_self_overlap}
    \lean{MPSTensor.isPrimitive_and_isNormal_of_irreducible_leftCanonical_selfOverlap_tendsto_one}
    \leanok
    \uses{def:irreducible_tensor, def:mpv_overlap, def:normal, thm:bnt_periodic_self_overlap_tendsto, thm:normal_tp_primitive_irred}
    Let $A$ be an irreducible left-canonical tensor of positive bond dimension.
    If $O_{AA}(N)\longrightarrow1$, then the transfer map of $A$ is primitive
    and $A$ is normal.
\end{theorem}

\begin{proof}
    \leanok
    \uses{thm:bnt_periodic_self_overlap_tendsto, thm:normal_tp_primitive_irred}
    Let $m$ be the peripheral period. Along its multiples,
    $O_{AA}(mn)\longrightarrow m$. The assumed full limit gives
    $O_{AA}(mn)\to1$, hence $m=1$. The transfer map is therefore primitive,
    and irreducibility and left-canonicality give normality.
\end{proof}

\begin{theorem}[Normality from normalized self-overlap]
    \label{thm:bnt_normality_from_self_overlap}
    \lean{MPSTensor.isNormal_of_irreducible_leftCanonical_selfOverlap_tendsto_one}
    \leanok
    \uses{thm:bnt_primitivity_normality_from_self_overlap}
    Let $A$ be an irreducible left-canonical tensor of positive bond dimension.
    If $O_{AA}(N)\to1$, then $A$ is normal.
\end{theorem}

\begin{proof}
    \leanok
    This is the normality conclusion of
    Theorem~\ref{thm:bnt_primitivity_normality_from_self_overlap}.
\end{proof}

\section{Primitive overlap convergence (complementary transfer-map gap form)}

The theorems in this section are stated directly in terms of the
complementary map $E-P$. This is the analytic form needed for the later
overlap argument, corresponding to the primitive-case specialization
of~\cite[Theorem~6.7]{Wolf2012Quantum} where $T_\phi$ is a rank-one
projection.
For primitive normalized transfer maps, the companion quantum-channel
volume~\cite[``Complementary transfer-map gap for primitive maps'']{QICLean2026Blueprint}
supplies this complementary transfer-map input under explicit hypotheses; the
rank-one Perron projection $P=P_\rho$, its trace, and the concrete
irreducible instantiation of the complementary spectral-radius gap are
given in Section~\ref{sec:app_spectral_perron_projection}.

\begin{theorem}[Complementary transfer-map gap implies trace convergence to $1$]%
    \label{thm:trace_pow_one}
    \lean{LinearMap.trace_pow_tendsto_one_of_spectralRadius_compl_lt_one}
    \leanok
    Let $E$ be trace-preserving with fixed point $\rho$, where
    $\tr(\rho)\neq0$, and let $P$ be the fixed-point projection of
    the corresponding definition in the companion quantum-channel volume~\cite[``Fixed-point projection'']{QICLean2026Blueprint}. If
    $\rho_{\spec}(E-P)<1$, then $\Tr(E^n)\to1$ as $n\to\infty$.
\end{theorem}

\begin{proof}
    \leanok

    By the power decomposition (established in the companion quantum-channel volume~\cite[``Power decomposition'']{QICLean2026Blueprint}),
    $E^n=P+(E-P)^n$ for $n\geq1$. Since $\rho_{\spec}(E-P)<1$, the
    Gelfand formula gives $(E-P)^n\to0$, so
    $\Tr(E^n)\to\Tr(P)$. Lemma~\cite[``Trace of the fixed-point projection'']{QICLean2026Blueprint}
    (Section~\ref{sec:app_spectral_perron_projection}) gives $\Tr(P)=1$.
\end{proof}

\begin{theorem}[Self-overlap convergence from a complementary transfer-map gap]%
    \label{thm:primitive_overlap}
    \lean{MPSTensor.mpvOverlap_tendsto_one_of_transfer_spectralRadius_compl_lt_one}
    \leanok
    \uses{def:transfer_map, def:mpv_overlap}
    Let $A$ be a normalized MPS tensor with a nonzero PSD fixed point
    $\rho$ of the transfer map. If $\rho_{\spec}(\E_A-P)<1$, where $P$ is
    the fixed-point projection, then $O_{AA}(N)\to1$ as $N\to\infty$.
\end{theorem}

\begin{proof}
    \leanok
    \uses{thm:trace_pow_one, thm:overlap_trace}
    By the overlap--trace identity (\ref{eq:spectral_overlap_trace}) with
    $A=B$, one has $O_{AA}(N)=\Tr(\E_A^N)$.
    By Theorem~\ref{thm:trace_pow_one}, $\Tr(\E_A^N)\to1$.
\end{proof}
